% !TeX root = ./../main.tex

\chapter{Prace powiązane z udziałem autora: struktura szumu kwantowego}\label{chapter:noise_structure}

\agentinfo{
Krótki rozdział (około 5 stron) omawiający dwie prace współautorskie, których pierwszym autorem jest Sebastian Raubitzek (SBA Research, Wiedeń).
Ustalenie autora z 2026-11-26: autor rozprawy implementował część kodu i uruchamiał eksperymenty.
Prace są przedstawiane jako prace powiązane z udziałem autora, a nie jako wkład tej rozprawy; rozdział nie ma własnej tezy.
Perspektywa uzupełniająca wobec rozdziałów 4--8: zamiast usuwać skutki szumu, opisuje się jego strukturę i możliwość jej wykorzystania.

Źródła:
rekonstrukcja szumu z generatorów łamiących symetrię, opublikowana w APS Open Science \agentref{bib:raubitzek_quantum_noise} \agentref{artykuły/quantum_noise_elsevier_submission/main.tex},
sieci neuronowe oparte na grupach Liego, preprint \agentref{bib:raubitzek_lie_group_manuscript} \agentref{artykuły/noisy_lie_group_learners/main.tex}.
Praca powiązana bez udziału autora rozprawy: raubitzek_obfuscation_2025.
Wersja w katalogu artykuły/ jest wersją zgłoszeniową; przed cytowaniem liczb sprawdzić zgodność z wersją opublikowaną.
Manuskrypt o grupach Liego jest bardzo obszerny; wybrać wyłącznie wątki związane z szumem i błędami.
Sformułowania muszą odpowiadać temu, co faktycznie pokazano: sieci są klasyczne, a zaburzenie łamiące symetrię jest opcjonalnym elementem konstrukcji.
Podstawy matematyczne są w zapleczu teoretycznym \agentref{chapters/theoretical_background_01_Q.tex} \agentref{chapters/theoretical_background_02_noise.tex}.
}

\section{Zakres i udział autora}\label{sec:ns:scope}
\agentinfo{
Jawne określenie udziału autora rozprawy w obu pracach (które części kodu, które eksperymenty), zgodne z oświadczeniami współautorów.
Uzasadnienie, dlaczego prace są omawiane: uzupełniają obraz szumu, z którym mierzą się metody z rozdziałów 4--8.
}

\section{Rekonstrukcja szumu z generatorów łamiących symetrię}\label{sec:ns:generators}
\agentinfo{
Bramka jako eksponenta generatora z addytywnym zaburzeniem; zaburzenie hermitowskie (błąd koherentny) i niehermitowskie (efektywna nieunitarność).
Analityczna jednokubitowa tomografia procesu, model próbkowania, dopasowanie zaburzeń algorytmem genetycznym.
Wynik: zaburzenia niehermitowskie odtwarzają macierz transferu Pauliego zmierzoną na urządzeniu IBM pełniej niż hermitowskie.
Zwięźle, bez powtarzania wyprowadzeń.
}

\section{Sieci oparte na grupach Liego i kontrolowane łamanie symetrii}\label{sec:ns:lie}
\agentinfo{
Warstwa działania grupy Liego: współrzędne w algebrze, mapa wykładnicza, lokalne działanie grupy.
Opcjonalne addytywne zaburzenie wzorowane na szumie kwantowym jako dodatkowy stopień swobody uczenia.
Sprawdzić w manuskrypcie, w których eksperymentach faktycznie użyto zaburzenia, i opisać tylko to, co wynika z danych.
}

\section{Znaczenie dla redukcji błędów}\label{sec:ns:implications}
\agentinfo{
Sygnatury łamania symetrii jako wskaźnik diagnostyczny oraz zastrzeżenia (stabilność zaburzeń, pominięte źródła szumu).
Wspólne założenie z metodą PDU: błąd jest stabilną, mierzalną cechą urządzenia.
Ograniczenia: model jednokubitowy, jedno urządzenie, model fenomenologiczny.
}
